% tranlational in lab methodology discussion 

Initial trials attempted to map the magnetic field from the furthest position of the pendulum, 8 cm from the solenoid face, at a constant solenoid current in a single run. The idea was to obtain ($B_0$) at the pendulum’s resting position and then compute the field gradient at any point from that map. Afterward, the pendulum was allowed to be slowly pulled toward the solenoid to its equilibrium position so its displacement (in cm) could be recorded. This was repeated while increasing the solenoid current in 0.5 A steps, aiming to fit the resulting equilibrium positions using Eq. \ref{eq:forFit}. 

These attempts produced inconsistent results that pointed to errors inherent to the setup. Because the current, rather than the entire apparatus, was changed, the value of $B_0$ at the pendulum’s resting position did not correspond to the required by Eq. \ref{eq:forFit}. Even using the mapped field, estimating the correct value for the magnetic field for each equilibrium position was unreliable. Moreover, predicting equilibrium positions before running the experiment proved difficult. Additional complications came from solenoid imperfections, current leakage, and induced currents, all of which became significant once displacements exceeded 1 cm. At higher fields the pendulum was sometimes pulled fully into the solenoid. Another facor was the gaussmeter restriction on accurate motion tracking to roughly 11–12 mm on the translation stage. While this approach resembles the proposed clinical protocol, the scale of the apparatus and the nonuniform fields make it impractical. Clinically, a single mapping session would give both B and $\nabla$B at any point, leaving the angle measurement as the critical step.


% discussion about translational clincal methodology


Initial setup revealed practical challenges with string attachment to the baskets. Knotting proved time-consuming and inconsistent across different basket designs, leading to variability in suspension length and pre-tension. To streamline future sessions, pre-prepared assemblies are recommended: additional 3D-printed baskets should be fabricated, each paired with a fixed-length string segment of approximately 1.3\,m, pre-weighted, and labeled for rapid deployment.

